\section{AMC 8 1987}
        \begin{enumerate}
        	\item (\textit{AMC 8 1987}) $.4+.02+.006=$


 $\text{(A)}\ .012 \qquad \text{(B)}\ .066 \qquad \text{(C)}\ .12 \qquad \text{(D)}\ .24 \qquad \text{(E)} .426$


 
 \textbf{Answer: }E



\textbf{Solution 1}

$.4+.02+.006 = .400+.020+.006 = \boxed{\textbf{(E)}\ .426}$

 



\textbf{Solution 2}

$.4 + .02 + .006 = \frac{4}{10} + \frac{2}{100} + \frac{6}{1000} = \frac{400}{1000} + \frac{20}{1000} + \frac{6}{1000} = \frac{426}{1000} = \boxed{\textbf{(E)}\ .426}$.

 


	\item (\textit{AMC 8 1987}) $\frac{2}{25}=$


 $\text{(A)}\ 0.008 \qquad \text{(B)}\ 0.08 \qquad \text{(C)}\ 0.8 \qquad \text{(D)} 1.25 \qquad \text{(E)}\ 12.5$


 
 \textbf{Answer: }B



\textbf{Solution}

$\frac{2}{25}=\frac{2\cdot 4}{25\cdot 4} = \frac{8}{100} = \boxed{\textbf{(B)}\ .08}$

 


	\item (\textit{AMC 8 1987}) $2(81+83+85+87+89+91+93+95+97+99)=$


 $\text{(A)}\ 1600 \qquad \text{(B)}\ 1650 \qquad \text{(C)}\ 1700 \qquad \text{(D)}\ 1750 \qquad \text{(E)}\ 1800$


 
 \textbf{Answer: }E



\textbf{Solution 1}

Rearranging terms yields $2(81 + 83 + 85 + 87 + 89 + 91 + 93 + 95 + 97 + 99)$ $= 2((81 + 99) + (83 + 97) + (85 + 95) + (87 + 93) + (89 + 91))$ 

 $= 2 (180 + 180 + 180 + 180 + 180)$ $= 2 \cdot 5 \cdot 180$ $= \boxed{\textbf{(E)}\ 1800}$.

 



\textbf{Solution 2}

Recall that the sum of an arithmetic series $a_1 + ... + a_n = \frac{n}{2}(a_1 + a_n)$, so this sum is $2\left(\frac{10}{2}(81 + 99)\right) = \boxed{\textbf{(E)}\ 1800}$.

 


	\item (\textit{AMC 8 1987}) Martians measure angles in clerts.  There are $500$ clerts in a full circle.  How many clerts are there in a right angle?


 $\text{(A)}\ 90 \qquad \text{(B)}\ 100 \qquad \text{(C)}\ 125 \qquad \text{(D)}\ 180 \qquad \text{(E)}\ 250$


 
 \textbf{Answer: }C



\textbf{Solution}

A right angle is $\frac{1}{4}$ of a circle, so it has $\frac{1}{4} \cdot 500=\boxed{\textbf{(C)}\ 125}$ clerts.

 


	\item (\textit{AMC 8 1987}) The area of the rectangular region is 


 
                    \begin{figure}[H]
                        \centering
                        \includegraphics[width=0.5\textwidth]{Problem/AMC_8_AJHSME_1987_5_Problem_Figure_1.png}
                    \end{figure}
                    


 $\text{(A)}\ \text{.088 m}^2 \qquad \text{(B)}\ \text{.62 m}^2 \qquad \text{(C)}\ \text{.88 m}^2 \qquad \text{(D)}\ \text{1.24 m}^2 \qquad \text{(E)}\ \text{4.22 m}^2$


 
 \textbf{Answer: }A



\textbf{Solution}

The area is $.4 \times .22 = \boxed{\textbf{(A)}\ 0.088}\ \text{m}^2$.

 


	\item (\textit{AMC 8 1987}) The smallest product one could obtain by multiplying two numbers in the set $\{ -7, -5, -1, 1, 3 \}$ is


 $\text{(A)}\ -35 \qquad \text{(B)}\ -21 \qquad \text{(C)}\ -15 \qquad \text{(D)}\ -1 \qquad \text{(E)}\ 3$


 
 \textbf{Answer: }B



\textbf{Solution}

Choosing the positive and negative numbers with the highest absolute values minimizes this product, which in this case is $-7 \cdot 3 = \boxed{\textbf{(B)}\ -21}$.

 


	\item (\textit{AMC 8 1987}) The large cube shown is made up of $27$ identical sized smaller cubes.  For each face of the large cube, the opposite face is shaded the same way.  The total number of smaller cubes that must have at least one face shaded is 


 $\text{(A)}\ 10 \qquad \text{(B)}\ 16 \qquad \text{(C)}\ 20 \qquad \text{(D)}\ 22 \qquad \text{(E)}\ 24$


 
                    \begin{figure}[H]
                        \centering
                        \includegraphics[width=0.5\textwidth]{Problem/AMC_8_AJHSME_1987_7_Problem_Figure_1.png}
                    \end{figure}
                    


 
 \textbf{Answer: }C



\textbf{Solution 1}

Notice that no two shaded tiles are adjacent.  Therefore, each small cube has at most one face shaded, and the number of small cubes with at least one face shaded is the same as the number of shaded tiles on the entire cube.  Since there are $10$ visible shaded tiles, and the hidden faces of the entire cube are identical to the shown faces, there must be $2 \cdot 10 = \boxed{\textbf{(C)}\ 20}$ small cubes with at least one face shaded.

 



\textbf{Solution 2}

The top and bottom layers each have $2$ cubes with no shaded sides, which are located in the middle of the front and back rows.  The middle layer has $3$ cubes with no shaded sides, which are all in the middle row.  Therefore, $2 \cdot 2 + 3 = 7$ cubes have no shaded sides, and $27 - 7 = \boxed{\textbf{(C)}\ 20}$ cubes have at least one shaded side.

 


	\item (\textit{AMC 8 1987}) If $\text{A}$ and $\text{B}$ are nonzero digits, then the number of digits (not necessarily different) in the sum of the three whole numbers is


 
 \[\begin{tabular}[t]{cccc} 9 & 8 & 7 & 6 \\ & A & 3 & 2 \\ & & B & 1 \\ \hline  \end{tabular}\]


 

$\text{(A)}\ 4 \qquad \text{(B)}\ 5 \qquad \text{(C)}\ 6 \qquad \text{(D)}\ 9 \qquad \text{(E)}\ \text{depends on the values of A and B}$


 
 \textbf{Answer: }B



\textbf{Solution}

The smallest possible sum occurs if $A=B=1$, so it is $9876+132+11=10019$.  The largest possible sum occurs if $A=B=9$, so it is $9876+932+91=10899$.  Therefore, the sum is between $10019$ and $10899$, so it has $\boxed{\textbf{(B)}\ 5}$ digits.

 


	\item (\textit{AMC 8 1987}) When finding the sum $\frac{1}{2}+\frac{1}{3}+\frac{1}{4}+\frac{1}{5}+\frac{1}{6}+\frac{1}{7}$, the least common denominator used is


 $\text{(A)}\ 120 \qquad \text{(B)}\ 210 \qquad \text{(C)}\ 420 \qquad \text{(D)}\ 840 \qquad \text{(E)}\ 5040$


 
 \textbf{Answer: }C



\textbf{Solution}

The least common denominator is $\mathrm{lcm}(2, 3, 4, 5, 6, 7) = \mathrm{lcm}(2, 3, 2^2, 5, 2\cdot 3, 7) = 2^2 \cdot 3 \cdot 5 \cdot 7 = \boxed{\textbf{(C)}\ 420}$.

 


	\item (\textit{AMC 8 1987}) $4(299)+3(299)+2(299)+298=$


 $\text{(A)}\ 2889 \qquad \text{(B)}\ 2989 \qquad \text{(C)}\ 2991 \qquad \text{(D)}\ 2999 \qquad \text{(E)}\ 3009$


 
 \textbf{Answer: }B



\textbf{Solution}

Using the distributive property, $4(299)+3(299)+2(299)+298$ $= 4(299)+3(299)+2(299)+1(299)-1$ $= (4+3+2+1)(299)-1$ $= 10(299)-1$ $= \boxed{\textbf{(B)}\ 2989}$.

 


	\item (\textit{AMC 8 1987}) The sum $2\frac17+3\frac12+5\frac{1}{19}$ is between


 $\text{(A)}\ 10\text{ and }10\frac12 \qquad \text{(B)}\ 10\frac12 \text{ and } 11 \qquad \text{(C)}\ 11\text{ and }11\frac12 \qquad \text{(D)}\ 11\frac12 \text{ and }12 \qquad \text{(E)}\ 12\text{ and }12\frac12$


 
 \textbf{Answer: }B



\textbf{Solution 1}

Since $\frac{1}{7}<\frac14$ and $\frac{1}{19}<\frac14$, it follows that $2\frac17+3\frac12+5\frac{1}{19}<2\frac14+3\frac12+5\frac14 =11$. Also, $2\frac17+3\frac12+5\frac{1}{19}>2+3\frac12+5=10\frac12$.  Therefore, the sum is between $\boxed{\textbf{(B)}\ 10\frac{1}{2}\text{ and }11}$.

 



\textbf{Solution 2}

This sum is $2 + \frac{1}{2} + 3 + \frac{1}{7} + 5 + \frac{1}{19} = 2 + 3 + 5 + \frac{1}{2} + \frac{1}{7} + \frac{1}{19} = 10\frac{1}{2} + \frac{1}{7} + \frac{1}{19} = 10\frac{1}{2} + \frac{26}{133}$.  Since $\frac{26}{133} < \frac{1}{2}$, this sum must be between $\boxed {\textbf {(B) } 10\frac12 \text{ and } 11}$.

 \href{https://artofproblemsolving.com/wiki/index.php?title=User:Anabel.disher}{Anabel.disher} (\href{https://artofproblemsolving.com/wiki/index.php?title=User\_talk:Anabel.disher}{talk}) 12:49, 18 June 2025 (EDT)

 



\textbf{Solution 3}

The total sum is $10 + \frac{1}{2} + \frac{1}{7} + \frac{1}{19} = 10 + \frac{9}{14} + \frac{1}{19} = 10 + \frac{185}{266} = 10 \frac{185}{266}$, which is between $\boxed {\textbf {(B) } 10\frac12 \text{ and } 11}$.

 


	\item (\textit{AMC 8 1987}) What fraction of the large $12$ by $18$ rectangular region is shaded?


 
                    \begin{figure}[H]
                        \centering
                        \includegraphics[width=0.5\textwidth]{Problem/AMC_8_AJHSME_1987_12_Problem_Figure_1.png}
                    \end{figure}
                    


 $\text{(A)}\ \frac{1}{108} \qquad \text{(B)}\ \frac{1}{18} \qquad \text{(C)}\ \frac{1}{12} \qquad \text{(D)}\ \frac29 \qquad \text{(E)}\ \frac13$


 
 \textbf{Answer: }C



\textbf{Solution 1}

The rectangle is divided into upper and lower halves, and then the lower half is divided into left and right halves.  The lower right section is therefore $\frac{1}{2} \cdot \frac{1}{2} = \frac{1}{4}$ of the area of the entire rectangle.  Since this section is divided into $6$ equal rectangles, $2$ of which are shaded, $\frac{1}{3}$ of the lower-right section is shaded, that is, $\frac{1}{3}\cdot \frac{1}{4}=\boxed{\textbf{(C)}\ \frac{1}{12}}$ of the large rectangle.

 



\textbf{Solution 2}

The total area of the rectangle is $12 \cdot 18 = 216$.  The $x$-coordinate of the left edge of the shaded area is $\frac{4}{6} = \frac{2}{3}$ of the way from $9$ to $18$, which is $9 + \frac{2}{3} \cdot 9 = 15$.  Therefore, the shaded area is $6 \cdot 3 = 18$.  Therefore, $\frac{18}{216} = \boxed{\textbf{(C)}\ \frac{1}{12}}$ of the entire rectangle is shaded.

 


	\item (\textit{AMC 8 1987}) Which of the following fractions has the largest value?


 $\text{(A)}\ \frac{3}{7} \qquad \text{(B)}\ \frac{4}{9} \qquad \text{(C)}\ \frac{17}{35} \qquad \text{(D)}\ \frac{100}{201} \qquad \text{(E)}\ \frac{151}{301}$


 
 \textbf{Answer: }E



\textbf{Solution}

Notice that $\frac{3}{7} < \frac{3}{6} = \frac{1}{2}$, $\frac{4}{9} < \frac{4}{8} = \frac{1}{2}$, $\frac{17}{35} < \frac{17}{34} = \frac{1}{2}$, and $\frac{100}{201} < \frac{100}{200} = \frac{1}{2}$, but $\frac{151}{301} > \frac{151}{302} = \frac{1}{2}$.  Since $\boxed{\textbf{(E)}\ \frac{151}{301}}$ is the only answer choice greater than $\frac{1}{2}$, it must be the largest.

 


	\item (\textit{AMC 8 1987}) A computer can do $10,000$ additions per second.  How many additions can it do in one hour?


 $\text{(A)}\ 6\text{ million} \qquad \text{(B)}\ 36\text{ million} \qquad \text{(C)}\ 60\text{ million} \qquad \text{(D)}\ 216\text{ million} \qquad \text{(E)}\ 360\text{ million}$


 
 \textbf{Answer: }B



\textbf{Solution}

Using dimensional analysis, ${\frac{10000\text{ additions}}{1\text{ second}} \cdot \frac{3600\text{ seconds}}{1\text{ hour}} = \boxed{\textbf{(B)}\ 36 \text{ million}}\text{ additions/hour}}$.

 


	\item (\textit{AMC 8 1987}) The sale ad read: "Buy three tires at the regular price and get the fourth tire for $\$3.$" Sam paid $\$240$ for a set of four tires at the sale.  What was the regular price of one tire?


 $\text{(A)}\ \$59.25 \qquad \text{(B)}\ \$60 \qquad \text{(C)}\ \$70 \qquad \text{(D)}\ \$79 \qquad \text{(E)}\ \$80$


 
 \textbf{Answer: }D



\textbf{Solution 1}

Let $x$ be the regular price of a tire.  Then $3x + \$3 = \$240$, so $x = \boxed{\textbf{(D)}\ \$79}$.

 



\textbf{Solution 2}

Since Sam paid $\$240$ for all $4$ tires, which included $\$3$ for the last one, he paid $\$240 - \$3 = \$237$ for the first $3$ tires at regular price.  Therefore, the regular price of a tire is $\$237 \div 3 = \boxed {\textbf {(D) } \$79}$.

 


	\item (\textit{AMC 8 1987}) Joyce made $12$ of her first $30$ shots in the first three games of this basketball season, so her seasonal shooting average was $40\%$.  In her next game, she took $10$ shots and raised her seasonal shooting average to $50\%$.  How many of these $10$ shots did she make?


 $\text{(A)}\ 2 \qquad \text{(B)}\ 3 \qquad \text{(C)}\ 5 \qquad \text{(D)}\ 6 \qquad \text{(E)}\ 8$


 
 \textbf{Answer: }E



\textbf{Solution}

After the fourth game, Joyce had taken $40$ total shots, $50\%$ of which she made, so she made $40\times .5=20$ shots.  Since she made $12$ shots in the first three games, she made $20-12=\boxed{\textbf{(E)}\ 8}$ shots in the last game.

 


	\item (\textit{AMC 8 1987}) Abby, Bret, Carl, and Dana are seated in a row of four seats numbered \#1 to \#4.  Joe looks at them and says:


 <pre>"Bret is next to Carl."
"Abby is between Bret and Carl."
</pre>However each one of Joe's statements is false.  Bret is actually sitting in seat \#3.  Who is sitting in seat \#2?


 $\text{(A)}\ \text{Abby} \qquad \text{(B)}\ \text{Bret} \qquad \text{(C)}\ \text{Carl} \qquad \text{(D)}\ \text{Dana} \qquad \text{(E)}\ \text{There is not enough information to be sure.}$


 
 \textbf{Answer: }D



\textbf{Solution}

Bret sits in seat $\#3$.

 
 \[\text{-----}\qquad \text{-----} \qquad \textbf{Bret} \qquad \text{-----}\]

 Since Carl is not next to Bret, he sits in seat $\#1$.

 
 \[\textbf{Carl}\qquad \text{-----} \qquad \text{Bret} \qquad \text{-----}\]

 Since Abby is not between Bret and Carl, she sits in seat $\#4$.  

 
 \[\text{Carl}\qquad \text{-----} \qquad \text{Bret} \qquad \textbf{Abby}\]

 This leaves $\boxed{\textbf{(D)}\ \text{Dana}}$ to sit in seat $\#2$.

 
 \[\text{Carl}\qquad \textbf{Dana} \qquad \text{Bret} \qquad \text{Abby}\]

 


	\item (\textit{AMC 8 1987}) Half the people in a room left.  One third of those remaining started to dance.  There were then $12$ people who were not dancing.  The original number of people in the room was what?


 $\text{(A)}\ 24 \qquad \text{(B)}\ 30 \qquad \text{(C)}\ 36 \qquad \text{(D)}\ 42 \qquad \text{(E)}\ 72$


 
 \textbf{Answer: }C



\textbf{Solution}

Let the original number of people in the room be $x$.  Half of them left, so $\frac{x}{2}$ of them are left in the room.

 After that, one third of this group is dancing, so $\frac{x}{2}-\frac{1}{3}\left( \frac{x}{2}\right) =\frac{x}{3}$ people are not dancing.

 This is given to be $12$, so 
 \[\frac{x}{3}=12\Rightarrow x=36\] 

 $\boxed{\text{C}}$

 


 



\textbf{Solution #2}

First note that of the $\frac{1}{2}$ people remaining in the room, $\frac{2}{3}$ are not dancing. Therefore $\frac{1}{2}\cdot\frac{2}{3}= \frac{1}{3}$ of the original amount of people in the room is $12$. The answer is $\boxed{C}$.

 


	\item (\textit{AMC 8 1987}) A calculator has a squaring key $\boxed{x^2}$ which replaces the current number displayed with its square.  For example, if the display is $\boxed{000003}$ and the $\boxed{x^2}$ key is depressed, then the display becomes $\boxed{000009}$.  If the display reads $\boxed{000002}$, how many times must you depress the $\boxed{x^2}$ key to produce a displayed number greater than $500$?


 $\text{(A)}\ 4 \qquad \text{(B)}\ 5 \qquad \text{(C)}\ 8 \qquad \text{(D)}\ 9 \qquad \text{(E)}\ 250$


 
 \textbf{Answer: }A



\textbf{Solution}

We just brute force this:
 \begin{align*} 2 &\rightarrow 4 \\ &\rightarrow 16 \\ &\rightarrow 256 \\ &\rightarrow 65536>500  \end{align*}


 Clearly we need to press the button $4$ times, so $\boxed{\text{A}}$

 


	\item (\textit{AMC 8 1987}) "If a whole number $n$ is not prime, then the whole number $n-2$ is not prime." A value of $n$ which shows this statement to be false is


 $\text{(A)}\ 9 \qquad \text{(B)}\ 12 \qquad \text{(C)}\ 13 \qquad \text{(D)}\ 16 \qquad \text{(E)}\ 23$


 
 \textbf{Answer: }A



\textbf{Solution}

To show that the conditional statement "If $P$, then $Q$" is false, we must find a counterexample where the starting condition ($P$) is true, but the conclusion ($Q$) is false. 

 In this problem:

 
                \begin{itemize}
                    \item Condition $P$: $n$ is not prime.\item Conclusion $Q$: $n-2$ is not prime.
                \end{itemize}


                Therefore, a counterexample requires a number $n$ that is \textbf{not prime} but results in an $n-2$ that \textbf{is prime} (making the conclusion false).

 Let's test the choices:

 
                \begin{itemize}
                    \item $\textbf{(A)}\ 9$: $9$ is not prime, and $9-2=7$, which \textbf{is} prime. This breaks the rule, so $9$ is a valid counterexample.\item $\textbf{(B)}\ 12$: $12$ is not prime, and $12-2=10$, which is not prime. This follows the rule.\item $\textbf{(C)}\ 13$: $13$ is prime, so it fails the initial condition.\item $\textbf{(D)}\ 16$: $16$ is not prime, and $16-2=14$, which is not prime. This follows the rule.\item $\textbf{(E)}\ 23$: $23$ is prime, so it fails the initial condition.
                \end{itemize}


                Only $\boxed{\textbf{(A)}\ 9}$ is a counterexample.

 


	\item (\textit{AMC 8 1987}) Suppose $n^{*}$ means $\frac{1}{n}$, the reciprocal of $n$.  For example, $5^{*}=\frac{1}{5}$.  How many of the following statements are true?


 <pre>i) $3^*+6^*=9^*$
ii) $6^*-4^*=2^*$
iii) $2^*\cdot 6^*=12^*$
iv) $10^*\div 2^* =5^*$
</pre>$\text{(A)}\ 0 \qquad \text{(B)}\ 1 \qquad \text{(C)}\ 2 \qquad \text{(D)}\ 3 \qquad \text{(E)}\ 4$


 
 \textbf{Answer: }C



\textbf{Solution}

Statement (i) is false, since $\frac{1}{3}+\frac{1}{6} = \frac{1}{2} \neq \frac{1}{9}$.

 Statement (ii) is false, since $\frac{1}{6}-\frac{1}{4} = -\frac{1}{12} \neq \frac{1}{2}$.

 Statement (iii) is true, since $\frac{1}{2}\cdot \frac{1}{6} = \frac{1}{12}$.

 Statement (iv) is true, since $\frac{1}{10}\div \frac{1}{2} = \frac{1}{5}$. 

 So $\boxed{\textbf{(C)}\ 2}$ of these statements are true.

 


	\item (\textit{AMC 8 1987}) $\text{ABCD}$ is a rectangle, $\text{D}$ is the center of the circle, and $\text{B}$ is on the circle.  If $\text{AD}=4$ and $\text{CD}=3$, then the area of the shaded region is between


 
                    \begin{figure}[H]
                        \centering
                        \includegraphics[width=0.5\textwidth]{Problem/AMC_8_AJHSME_1987_22_Problem_Figure_1.png}
                    \end{figure}
                    


 $\text{(A)}\ 4\text{ and }5 \qquad \text{(B)}\ 5\text{ and }6 \qquad \text{(C)}\ 6\text{ and }7 \qquad \text{(D)}\ 7\text{ and }8 \qquad \text{(E)}\ 8\text{ and }9$


 
 \textbf{Answer: }D



\textbf{Solution}

Since $ABCD$ is a rectangle, $BC = AD = 4$.  So by the Pythagorean theorem, $BD = \sqrt{BC^2 + CD^2} = \sqrt{3^2+4^2} = 5$.  So, the circle has radius $5$ and total area $25\pi$.  Consequently, the quarter circle bounded by $\angle D$ has area $\frac{25\pi}{4}$.  

 Since $ABCD$ has area $3\cdot 4=12$, the shaded area is $\frac{25\pi}{4}-12 \approx \frac{25 \cdot 3.14}{4} - 12 = 7.625$, so the area is between $\boxed{\textbf{(D)}\ \text{7 and 8}}$.

 


	\item (\textit{AMC 8 1987}) Assume the adjoining chart shows the $1980$ U.S. population, in millions, for each region by ethnic group.  To the nearest percent, what percent of the U.S. Black population lived in the South?


 
 \[\begin{tabular}[t]{c|cccc}  & NE & MW & South & West \\ \hline White & 42 & 52 & 57 & 35 \\ Black & 5 & 5 & 15 & 2 \\ Asian & 1 & 1 & 1 & 3 \\ Other & 1 & 1 & 2 & 4  \end{tabular}\]


 $\text{(A)}\ 20\% \qquad \text{(B)}\ 25\% \qquad \text{(C)}\ 40\% \qquad \text{(D)}\ 56\% \qquad \text{(E)}\ 80\%$


 
 \textbf{Answer: }D



\textbf{Solution}

The total U.S. Black population was $5+5+15+2=27$ million.  So $\frac{15}{27}= \frac{5}{9} \approx \boxed{\textbf{(D)}\ 56\%}$ of the U.S. Black population lived in the South.

 


	\item (\textit{AMC 8 1987}) A multiple choice examination consists of $20$ questions.  The scoring is $+5$ for each correct answer, $-2$ for each incorrect answer, and $0$ for each unanswered question.  John's score on the examination is $48$.  What is the maximum number of questions he could have answered correctly?


 $\text{(A)}\ 9 \qquad \text{(B)}\ 10 \qquad \text{(C)}\ 11 \qquad \text{(D)}\ 12 \qquad \text{(E)}\ 16$


 
 \textbf{Answer: }D



\textbf{Solution 1}

Let $x$ be the number of questions John answered correctly.  Since he answered at most $20-x$ questions incorrectly, he scored at least $5x - 2(20-x) = 7x - 40$ points.  So $7x - 40 \leq 48$, and thus $x \leq \frac{88}{7} = 12\frac{4}{7}$.  Therefore, he answered at most $\boxed{\textbf{(D)}\  12}$ questions correctly, which would result in a score of $48$ if he answered $6$ questions incorrectly.

 



\textbf{Solution 2}

Since John scored $48$ points, and each correct answer scores $5$ points, he answered at least $10$ questions correctly.

 If he answered $10$ questions correctly, he gained $50$ points from correct answers, so he lost $2$ points from incorrect answers.  Therefore, he could have scored $48$ points by answering $10$ questions correctly and $1$ question incorrectly.

 If he answered $11$ questions correctly, he gained $55$ points from correct answers, so he lost $7$ points from incorrect answers.  This is impossible since $7$ is not divisible by $2$.

 If he answered $12$ questions correctly, he gained $60$ points from correct answers, so he lost $12$ points from incorrect answers.  Therefore, he could have scored $48$ points by answering $12$ questions correctly and $6$ questions incorrectly.

 If he answered $13$ or more questions correctly, he gained at least $65$ points from correct answers, so he lost at least $17$ points from incorrect answers.  So, he answered at least $9$ questions incorrectly, which would be impossible since the test only has $20$ questions.

 Therefore, he answered at most $\boxed{\textbf{(D)}\ 12}$ questions correctly.

 


 



\textbf{Solution 3 (Answer Choices)}

If John answered $16$ questions correctly, he gained $80$ points from correct answers, so he lost $32$ points from incorrect answers.  So, he answered $16$ questions incorrectly, which is impossible since the test only had $20$ questions.

 If he answered $12$ questions correctly, he gained $60$ points from correct answers, so he lost $12$ points from incorrect answers.  So, he answered $6$ questions incorrectly, which is possible.  

 Since all the other answer choices are lower than $12$, he answered at most $\boxed{\textbf{(D)}\ 12}$ questions correctly.

 


	\item (\textit{AMC 8 1987}) Ten balls numbered $1$ to $10$ are in a jar.  Jack reaches into the jar and randomly removes one of the balls.  Then Jill reaches into the jar and randomly removes a different ball.  The probability that the sum of the two numbers on the balls removed is even is


 $\text{(A)}\ \frac{4}{9} \qquad \text{(B)}\ \frac{9}{19} \qquad \text{(C)}\ \frac{1}{2} \qquad \text{(D)}\ \frac{10}{19} \qquad \text{(E)}\ \frac{5}{9}$


 
 \textbf{Answer: }A



\textbf{Solution 1}

At first, there are $5$ even balls and $5$ odd balls.  Therefore, after Jack chooses a ball, the jar will have $4$ balls with the same parity (odd/even) as Jack's ball and $5$ balls with the opposite parity.  For the sum to be even, Jill must choose a ball with the same parity as Jack's ball.  The probability that this happens is $\boxed{\textbf{(A)}\ \frac{4}{9}}$.

 



\textbf{Solution 2}

Jack has $10$ choices for his ball, and after he makes his selection Jill will have $9$ choices for hers.  Therefore, there are $10 \cdot 9 = 90$ combinations of selections, and each combination is equally likely.

 For the sum of their numbers to be even, Jack and Jill must either both select odd numbers or both select even numbers.  

 If they both select even numbers, Jack has $5$ choices for his ball, and Jill will have $4$ choices for hers, for a total of $4\cdot5 = 20$ combinations of selections.

 If they both select odd numbers, Jack has $5$ choices for his ball, and Jill will have $4$ choices for hers, for a total of $4\cdot5 = 20$ combinations of selections.

 Therefore, of the $90$ combinations of selections, $20+20=40$ yield an even sum.  Therefore, the probability that the sum is even is $\boxed{\textbf{(A)}\ \frac{4}{9}}$.

 



\textbf{Solution 3 - Complementary Counting}

An alternative approach would be to calculate the probability the sum of the two numbers is odd, and subtract that from the total. We see that there is only one way to get an odd sum, where the two numbers must be of opposite parity.

 Here, we have two cases $\rightarrow$ the two numbers are, respectively, (even, odd), and (odd, even). The probability to get an even number as the first number is $\frac{5}{10} = \frac{1}{2},$ and an odd number as the second number gives us a probability of $\frac{5}{9}.$ Multiplying the two, we get $\frac{5}{18},$ and because both cases have equal probability, we have the probability for the sum to be odd is $\frac{10}{18} = \frac{5}{9}.$

 Finally, subtracting this from $1$ gives us $\boxed{\frac{4}{9} \text{ } (A)}.$~ Cheetahboy93

 


\end{enumerate}